\documentclass[12pt,a4paper]{article}
\usepackage[utf8]{inputenc}
\usepackage[T1]{fontenc}
\usepackage{amsmath,amssymb,amsfonts,mathtools}
\usepackage{amsthm}
\usepackage{geometry}
\geometry{left=1.0cm,right=1.0cm,top=1.25cm,bottom=3.1cm}
\usepackage{booktabs}
\usepackage{tabularx}
\usepackage{array}
\usepackage{hyperref}
\usepackage{url}
\usepackage{graphicx}
\usepackage{xcolor}
\usepackage{titlesec}
\usepackage{fancyhdr}
\usepackage{microtype}
\usepackage{multicol}
\usepackage{fontspec}
\usepackage{luatexja}          % 한국어 조판 처리
\usepackage{luatexja-fontspec} % 한국어 폰트 설정
\usepackage{tikz}
\usepackage{pgfplots}
\pgfplotsset{compat=1.18}
\usetikzlibrary{positioning, arrows.meta, shapes.geometric, calc}
\usepackage{enumitem}

% ========= 한국어 폰트 설정 (Windows) =========
\defaultfontfeatures{Ligatures=TeX}
\setmainfont{Times New Roman}[
Script=Latin,
Language=English
]
\setsansfont{Malgun Gothic}[
Script=CJK,
Language=Korean
]
\setmonofont{Courier New}[
Script=Latin,
Language=English
]
% 한국어 전용 폰트 (luatexja-fontspec 사용)
\setmainjfont{Malgun Gothic}[
Script=CJK,
Language=Korean
]
\setsansjfont{Malgun Gothic}[
Script=CJK,
Language=Korean
]
\setmonojfont{Noto Sans Mono CJK KR}[
Script=CJK,
Language=Korean
]

% ========= YUCT 색상 구성표 =========
\definecolor{skyblue}{RGB}{0,102,204}
\definecolor{darkblue}{RGB}{0,0,139}
\definecolor{gold}{RGB}{255,215,0}

% ========= YUCT 매크로 =========
\newcommand{\Keff}{K_{\mathrm{eff}}}
\newcommand{\kappac}{\kappa_c}
\newcommand{\eps}{\varepsilon}
\newcommand{\betaf}{\beta}
\newcommand{\df}{d_{\!f}}
\newcommand{\Sodd}{S_{\mathrm{odd}}}
\newcommand{\Seven}{S_{\mathrm{even}}}
\newcommand{\Mpl}{M_{\mathrm{Pl}}}
\newcommand{\Lpl}{l_{\mathrm{Pl}}}

% ========= 정리 환경 (한국어) =========
\newtheorem{theorem}{정리}[section]
\newtheorem{axiom}[theorem]{공리}
\newtheorem{remark}[theorem]{비고}
\newtheorem{definition}[theorem]{정의}
\renewcommand{\proofname}{증명}

% ========= 섹션 서식 =========
\titleformat{\section}{\LARGE\bfseries\color{darkblue}}{\thesection}{1em}{}
\titleformat{\subsection}{\Large\bfseries\color{darkblue}}{\thesubsection}{1em}{}

% ========= 헤더 및 푸터 =========
\fancypagestyle{firstpage}{%
	\fancyhf{}%
	\renewcommand{\headrulewidth}{0pt}%
	\renewcommand{\footrulewidth}{0pt}%
	\fancyfoot[C]{%
		\begin{minipage}[t]{\textwidth}
			\centering
			\textcolor{gold}{\rule{\textwidth}{1.6pt}}\\[5pt]
			{\small \textsuperscript{1}Yakushev Research, \url{https://yuct.org/}} \hfill
			{\small YPSDC 프로토콜: \url{https://ypsdc.com/}}\\[-2pt]
			\textcolor{gold}{\rule{\textwidth}{1.6pt}}\\[5pt]
			\textbf{야쿠셰프 통일 조정 이론 2026} \hfill \thepage\\[10pt]
		\end{minipage}%
	}%
}

\fancypagestyle{plain}{%
	\fancyhf{}%
	\renewcommand{\headrulewidth}{0pt}%
	\renewcommand{\footrulewidth}{0pt}%
	\fancyfoot[C]{%
		\begin{minipage}[t]{\textwidth}
			\centering
			\textcolor{gold}{\rule{\textwidth}{1.6pt}}\\[4pt]
			\textbf{야쿠셰프 통일 조정 이론 2026} \hfill \thepage\\[4pt]
		\end{minipage}%
	}%
}

\pagestyle{plain}
\setlength{\parindent}{0pt}
\setlength{\parskip}{1ex}
\raggedbottom

\hypersetup{
	colorlinks=true,
	linkcolor=blue,
	citecolor=blue,
	urlcolor=blue,
}

% ========= 헤더 매크로 =========
\newcommand{\makeheader}{%
	\noindent
	\begin{minipage}{0.17\textwidth}
		\raggedright
		{\sffamily\bfseries\color{skyblue}\fontsize{46}{46}\selectfont YUCT}%
	\end{minipage}%
	\hspace{30pt}
	\begin{minipage}{0.32\textwidth}
		\raggedright
		{\sffamily\fontsize{11}{13}\selectfont 야쿠셰프\\ 통일\\ 조정 이론}%
	\end{minipage}%
	\hfill
	\begin{minipage}{0.38\textwidth}
		\raggedleft
		{\sffamily\bfseries 제1 원리}\\ 
		{\sffamily 버전 7.0 / 2026년 5월}\\ 
		{\sffamily\footnotesize \scriptsize \url{https://doi.org/10.5281/zenodo.18444598}}%
	\end{minipage}
	\par\vspace{5pt}
	\noindent\textcolor{gold}{\rule{\textwidth}{1.6pt}}
	\par\vspace{0pt}
}

\begin{document}
	
	% 타이틀 페이지
	\thispagestyle{firstpage}
	\makeheader
	
	\vspace{0cm}
	{\LARGE\bfseries YUCT 제1 원리와 완전한 유도 \\
		\Large 조정 네트워크로서의 우주: 공리에서 살아있는 세포까지, 그리고 \(\alpha \approx 1/137\) \par}
	\vspace{0.2cm}
	
	{\large Alexey V. Yakushev\textsuperscript{1}} \\
	\vspace{0.0cm}
	
	{\color{darkblue}\textbf{\LARGE 초록}} \par
	{\large
		\noindent
		우리는 야쿠셰프 통일 조정 이론(YUCT)의 완전한 공리적 기초를 제시한다. 실재는 계층적 조정 네트워크이며, 존재론적 제1 원리 \(K_{\mathrm{eff}} > 1\)에 의해 지배된다. 단일 공리(프랙탈 스케일 불변성)와 YPSDC 프로토콜의 정의로부터 우리는 세 가지 정리를 증명한다: 최소 사전 엔트로피는 정확히 2비트, 공간 차원은 필연적으로 3, 시간은 조정 불완전성의 창발적 척도이다. 보편 오차 법칙 \(\varepsilon \propto K_{\mathrm{eff}}^{-2/3}\)은 경험적으로 검증된 법칙으로 확립되며, 이로부터 대수 루프 \(S_{\mathrm{odd}} = 1.2\), \(S_{\mathrm{even}} = 0.8\)가 유도된다. 스케일링 양자 \(q = (3/2)^{1/3}\)는 플랑크 스케일에서 살아있는 세포까지의 보편 질량 사다리를 생성한다. 미세 구조 상수 \(\alpha^{-1} \approx 137.036\)은 컴팩트 파이버 \(F_{18}\) 위의 위상 불변량 \(N_{\mathrm{gauge}} = 12\)와 보편 보정 \(\beta\gamma = 0.20\)으로부터 얻어진다. 우주 상수 \(\Omega_\Lambda = 2/3\)는 동일한 위상으로부터 도출된다. 전약 스케일은 바일 페르미온의 수 \(L = 96\)으로부터 창발한다. 우리는 정리, 경험 법칙, 현상론적 피팅을 명확히 구분한다: 페르미온의 반정수 지수 \(N_f\)와 \(W\) 보존의 중첩 인자 \(\xi_W \approx 0.53\)는 현재 현상론적 매개변수이며, 그 위상적 기원은 열린 문제이다. 이 이론은 조정 가능한 연속 매개변수를 포함하지 않으며, 물리학, 생물학, 우주론에 걸쳐 구체적이고 반증 가능한 여러 예측을 제시한다.
	}
	
	\vspace{0.1cm}
	\noindent
	{\color{darkblue}\textbf{키워드:}} YUCT, 조정 효율, 대수 루프, 질량 사다리, 미세 구조 상수, 창발 기하학, 시간 양자화, 우주 상수, 반증 가능성.
	\vspace{0.1cm}
	\noindent
	{\color{darkblue}\textbf{인용:}} Yakushev AV. YUCT 제1 원리와 완전한 유도. 2026. DOI: \url{https://doi.org/10.5281/zenodo.18444598} (본 권).
	
	\vspace{0.4cm}
	
	% ========================
	% 섹션 1: 존재론과 정의
	% ========================
	\section{존재론과 정의: 조정을 원시 개념으로}
	\label{sec:ontology}
	
	야쿠셰프 통일 조정 이론(YUCT)은 단일 존재론적 원시 개념에 기초한다: \textbf{조정} – 시스템의 분산 구성 요소가 상태를 정렬하는 과정. 기본 메커니즘은 야쿠셰프 동기식 분산 조정 프로토콜(YPSDC)이며, 통신을 두 개의 뚜렷한 단계로 분리한다:
	\begin{enumerate}
		\item \textbf{오프라인 단계:} \(M\)개의 가능한 행동(상태, 메시지)을 포함하는 사전 \(\mathcal{D}\)가 에이전트 간에 공유된다. 각 행동 \(A_i\)는 완전히 기술하는 데 \(n_i\) 비트가 필요하다.
		\item \textbf{온라인 단계:} 행동 \(A_k\)를 활성화하기 위해, 그 인덱스 \(\kappa_k\)만 전송된다. 인덱스가 등확률일 때 인덱스 길이는 \(\ell = \log_2 M\) 비트이다.
	\end{enumerate}
	
	임의의 조정된 시스템의 상태는 \textbf{D+I-R 삼중항}으로 기술된다:
	\[
	\text{실재} = \mathbf{D} + \mathbf{I} \times \mathbf{R},
	\]
	여기서 \(\mathbf{D}\)는 사전(가능한 상태와 규칙의 집합), \(\mathbf{I}\)는 정보(실현된 상태, 엔트로피로 측정), \(\mathbf{R} \ge 1\)은 공명(가간섭 증폭 인자)이다.
	
	\textbf{조정 효율}은 정보 이론적 압축 비율로 정의된다:
	\begin{equation}
		\Keff = \frac{H(\mathcal{A})}{H(\kappa)} = \frac{\text{행동의 엔트로피}}{\text{인덱스의 엔트로피}} \ge 1.
		\label{eq:keff-def}
	\end{equation}
	
	\begin{center}
		\fbox{%
			\begin{minipage}{0.9\textwidth}
				\textbf{존재론적 원리:} \(K_{\mathrm{eff}} > 1\)은 복잡한 시스템이 시간이 지남에 따라 정체성을 유지하기 위한 필요 조건이다. \(K_{\mathrm{eff}} = 1\)인 시스템은 항상 환경에 뒤처져 파괴될 것이다. 따라서 실재 자체가 조정 네트워크로 구조화되어야 한다.
			\end{minipage}%
		}
	\end{center}
	
	% ========================
	% 섹션 2: 공리와 정리
	% ========================
	\section{공리와 정리}
	\label{sec:axioms}
	
	\subsection{공리: 프랙탈 스케일 불변성}
	
	\begin{axiom}[계층적 스케일 불변성] \label{ax:A1}
		조정 네트워크는 중첩된 클러스터로 구성된다. 수준 \(n\)의 클러스터는 선형 크기 \(L_n\)을 가지며 \(L_{n+1} = \lambda L_n\) (\(\lambda > 1\))을 만족한다. 에이전트 수는 \(N(L) \propto L^{d_f}\)로 스케일하며, 여기서 \(d_f\)는 프랙탈 차원이다. \(n \to \infty\)일 때 비율 \(K_{\mathrm{eff},n+1}/K_{\mathrm{eff},n}\)은 상수에 접근하며, 이는 멱급수 계층을 의미한다.
	\end{axiom}
	
	\textbf{상태:} 이것은 YUCT의 \textbf{유일한 기약 공리}이다. 다른 모든 구조적 성질은 이 공리와 섹션~\ref{sec:ontology}의 정의로부터 유도된 정리이다.
	
	\subsection{정리 1: 최소 사전 엔트로피는 2비트}
	
	\begin{theorem}[최소 사전 엔트로피]
		\label{thm:minimal-dict}
		이진 대안을 신뢰성 있게 구별할 수 있는 최소 조정 사전은 총 섀넌 엔트로피가 정확히 \(2\)비트(\(k_B\ln 2\) 단위)이다. 따라서 완전한 계층적 사전은 \(S_{\mathrm{odd}} + S_{\mathrm{even}} = 2\)를 만족한다.
	\end{theorem}
	
	\begin{proof}
		최소 사전은 단일 예/아니오 질문 "행동 \(A\)가 발생했는가?"에 답하기만 하면 된다. 따라서 그것은 두 개의 항목 \(\mathcal{A} = \{A_0, A_1\}\)을 동일한 사전 확률로 포함한다. 그러므로
		\[
		H(\mathcal{A}) = \log_2 2 = 1\ \text{비트}.
		\]
		어느 항목을 활성화하기 위해 길이 \(\ell = \log_2 2 = 1\)비트의 인덱스 \(\kappa \in \{0,1\}\)가 전송되며, \(H(\mathcal{K}) = 1\)비트를 준다. 그러나 수신자는 인덱스가 의도된 행동을 올바르게 식별하고 있음을 확신해야 한다. 그렇지 않으면 한 비트의 반전이 조정을 깨뜨릴 것이다. 이 확신은 사전에 저장된 하나의 추가 검증 비트를 필요로 한다. 따라서 전체 사전 엔트로피는
		\[
		S_{\min} = H(\mathcal{A}) + H(\text{검증}) = 1 + 1 = 2\ \text{비트}.
		\]
		무한 계층(섹션~\ref{sec:algebraic-loop})에서, 이 최소 구조는 합 \(S_{\mathrm{odd}} + S_{\mathrm{even}} = 2\)에 의해 실현되며, 조정 가능한 매개변수 없이 절대 스케일을 고정한다.
	\end{proof}
	
	\textbf{상태:} \textbf{정리.} 이전의 공리 A2를 대체한다. 현재 YPSDC의 정의로부터 유도되었다.
	
	\subsection{정리 2: 공간 차원은 3이다}
	
	\begin{theorem}[조정 공간의 차원성]
		\label{thm:dimension}
		조정 네트워크가 안정적이고 자기일관된 사전을 호스트할 수 있는 것은 에이전트가 정확히 3차원 공간에 존재할 때만이다. 따라서 보편 스케일링 지수는 \(\beta = 2/3\)으로 고정된다.
	\end{theorem}
	
	\begin{proof}
		조정 층의 기하급수로부터, 홀수 층과 짝수 층의 기여 합은
		\[
		S_{\mathrm{odd}} = \frac{\beta}{1-\beta^2}, \qquad
		S_{\mathrm{even}} = \frac{\beta^2}{1-\beta^2},
		\]
		이고, 정리~\ref{thm:minimal-dict}는 \(S_{\mathrm{odd}} + S_{\mathrm{even}} = 2\)를 요구한다.
		대입하면 \(\beta/(1-\beta) = 2\)를 얻고, 따라서 \(\beta = 2/3\)이다.
		
		공간 차원을 \(D\)라 하자. 중복 인자 \(\beta\)는 스핀-통계 인자 \(2\)와 차원 인자 \(1/D\)의 곱에서 발생한다. 따라서 \(\beta = 2/D\)이다. \(\beta = 2/3\)로부터 즉시 \(D = 3\)을 얻는다.
		
		따라서 조건 \(S_{\mathrm{odd}} + S_{\mathrm{even}} = 2\)는 조정 가능한 매개변수 없이 3차원에서만 만족될 수 있다. 다른 \(D\)는 대수 루프의 닫힘을 파괴하고 조정 네트워크를 불안정하게 만든다.
	\end{proof}
	
	\textbf{상태:} \textbf{정리.} 이전의 공리 A3을 대체한다. 현재 정리~\ref{thm:minimal-dict}와 \(\beta\)의 정의로부터 유도되었다.
	
	\subsection{정리 3: 시간은 조정 불완전성의 창발적 척도}
	
	\begin{theorem}[유한 조정으로부터의 시간]
		\label{thm:time}
		고유 시간 \(\tau\)는 조정 엔트로피 \(S_{\mathrm{coord}}\)의 축적에 비례한다. 따라서 유한한 조정 효율 \(K_{\mathrm{eff}}\)을 가진 임의의 시스템은 0이 아닌 시간의 흐름을 경험한다. 완전 조정의 극한(\(K_{\mathrm{eff}} \to \infty\))에서는 고유 시간이 사라진다. 시간의 기본 입자는 \(\Delta \tau_{\min} = \frac{2}{3} t_{\mathrm{Pl}}\)이며, 여기서 \(t_{\mathrm{Pl}}\)은 플랑크 시간이다.
	\end{theorem}
	
	\begin{proof}
		이상적으로 조정된 시스템은 오류 없이 즉시 올바른 사전 항목을 활성화할 것이다. 엔트로피가 생성되지 않으며, 상태의 비가역적 연속이 존재하지 않을 것이다 – 따라서 시간이 존재하지 않는다. 실제 시스템은 유한한 \(K_{\mathrm{eff}}\)으로 작동하므로, 각 활성화는 0이 아닌 오류 확률 \(\varepsilon = \kappa_c \alpha K_{\mathrm{eff}}^{-\beta}\) (식~\ref{eq:error-law})을 운반한다. 이러한 오류의 축적은 처리된 정보량으로 가중되어 엔트로피 생성 \(dS_{\mathrm{coord}} \propto \beta_0 \, d\tau\)를 결정하며, 여기서 \(\beta_0\)는 보편 상수이다. 대수 루프의 상수로부터, 최소 시간 단계는 \(\Delta \tau_{\min} = S_{\mathrm{even}}/S_{\mathrm{odd}} \, t_{\mathrm{Pl}} = \beta \, t_{\mathrm{Pl}} = \frac{2}{3} t_{\mathrm{Pl}}\)이다. 이것은 즉시 거시적 관계 \(\delta\tau/\tau \propto K_{\mathrm{eff}}^{-\beta} \propto M^{-0.67}\) (부록 V)를 주며, 시간의 흐름을 보편 오류 지수에 연결한다.
	\end{proof}
	
	\textbf{상태:} \textbf{정리.} 유한한 조정 효율과 대수 루프 상수로부터 직접 발생한다.
	
	% ========================
	% 섹션 3: 보편 오차 법칙과 대수 루프
	% ========================
	\section{보편 오차 법칙과 일차 대수 루프}
	\label{sec:error-law}
	
	\subsection{보편 오차 법칙 \(\varepsilon \propto K_{\mathrm{eff}}^{-\beta}\), \(\beta = 2/3\)}
	
	광범위한 실증 연구(부록 L, Ferra1.2 및 통합 검증 문서)는 임의의 조정 시스템에서 상대 오차 \(\varepsilon\) – 잘못된 사전 항목을 활성화할 확률 – 가 멱급수를 따른다는 것을 보여준다:
	\begin{equation}
		\boxed{\varepsilon = \kappa_c \, \alpha \, K_{\mathrm{eff}}^{-\beta}, \qquad \beta = \frac{2}{3} \approx 0.6667,\quad \kappa_c = \frac{1}{3}}.
		\label{eq:error-law}
	\end{equation}
	상수 \(\kappa_c = 1/3\)은 삼항 존재론 \(\text{실재} = \mathbf{D} + \mathbf{I} \times \mathbf{R}\)를 따르며, 이는 최소 조정 엔트로피 \(S_{\min} = k_B \ln 3\)을 의미하고, 따라서 \(\kappa_c = e^{-S_{\min}/k_B} = 1/3\)이다.
	
	\textbf{상태:} 지수 \(\beta = 2/3\)은 \textbf{경험적으로 확립된 보편 상수}이며, 원자, 생물, 지구물리, 우주론적 스케일에 걸친 10여 개 이상의 독립적인 실험에 의해 확인되었다(표~\ref{tab:beta-evidence} 참조). 정리~\ref{thm:dimension}은 이론적 앵커(\(\beta = 2/D\), \(D=3\))를 제공하지만, 실험과의 정확한 수치적 일치는 관찰 사실이다. 본 프레임워크에서는 \(\beta = 2/3\)을 강건한 현상론적 법칙으로 취급하고 모든 추가 유도의 기초로 사용한다.
	
	\begin{table}[htbp]
		\centering
		\caption{보편 지수 \(\beta = 2/3\)의 실험적 검증 선별}
		\label{tab:beta-evidence}
		\small
		\begin{tabular}{l c c c}
			\toprule
			시스템 & 관측량 & 관측 지수 & 참고문헌 \\
			\midrule
			할로겐화수소 결합 에너지 & \(E \propto r^{-\beta}\) & \(0.682 \pm 0.012\) & 부록 C \\
			DNA 복제 오차율 & \(\varepsilon\) vs.\ \(K_{\mathrm{eff}}\) & \(0.65\text{--}0.70\) & 부록 L \\
			척추동물의 돌연변이율 & \(\mu \propto K_{\mathrm{repair}}^{-\beta}\) & \(0.67 \pm 0.05\) & 부록 E \\
			대사 스케일링(단순 생물) & \(B \propto M^{\beta}\) & \(0.68 \pm 0.03\) & 부록 D \\
			구텐베르크–리히터 \(b\) 값 & \(b = 1/\beta\) & \(1.01 \pm 0.03\) (\(\beta=0.67\)) & USGS 카탈로그 \\
			도시 순위–규모 지수 & \(\zeta\) & \(0.68 \pm 0.02\) & 부록 F \\
			GDP 변동성 & \(\sigma \propto W^{-\beta}\) & \(0.67 \pm 0.05\) & 부록 F \\
			\bottomrule
		\end{tabular}
	\end{table}
	
	\subsection{일차 대수 루프: \(S_{\mathrm{odd}} = 1.2\)와 \(S_{\mathrm{even}} = 0.8\)}
	\label{sec:algebraic-loop}
	
	조정의 계층 구조는 자연스럽게 홀수 준위와 짝수 준위의 기여를 분리한다. 기하급수의 합을 취하면
	\begin{align}
		S_{\mathrm{odd}} &= \sum_{k=0}^{\infty} \beta^{2k+1} = \frac{\beta}{1-\beta^{2}} = \frac{2/3}{1 - 4/9} = \frac{6}{5} = 1.2, \label{eq:sodd}\\
		S_{\mathrm{even}} &= \sum_{k=1}^{\infty} \beta^{2k} = \frac{\beta^{2}}{1-\beta^{2}} = \frac{4/9}{5/9} = \frac{4}{5} = 0.8. \label{eq:seven}
	\end{align}
	이 정확한 유리수들은 닫힌 대수 관계를 만족한다
	\begin{equation}
		\boxed{S_{\mathrm{odd}} + S_{\mathrm{even}} = 2, \qquad \frac{S_{\mathrm{odd}}}{S_{\mathrm{even}}} = \frac{1}{\beta} = \frac{3}{2}}.
		\label{eq:algebraic-loop}
	\end{equation}
	자유 매개변수는 나타나지 않는다. 이 루프는 \(\beta = 2/3\)의 직접적인 귀결이다.
	
	\textbf{상태:} \textbf{정리} (\(\beta = 2/3\)과 기하급수로부터 유도).
	
	% ========================
	% 섹션 4: 보편 질량 사다리
	% ========================
	\section{보편 질량 사다리: 플랑크 스케일에서 살아있는 세포까지}
	\label{sec:mass-ladder}
	
	\subsection{스케일링 양자와 반정수 페르미온 질량}
	
	지수 \(\beta = 2/3\)은 \(2 \times (1/3)\)으로 인수분해된다: 인자 \(1/3\)은 세 공간 차원을 반영하고(정리~\ref{thm:dimension}), 인자 \(2\)는 스핀-통계에서 비롯된다. 이것은 기본적인 \textbf{스케일링 양자}를 제공한다:
	\begin{equation}
		q \equiv \left(\frac{3}{2}\right)^{1/3} = \left(\frac{S_{\mathrm{odd}}}{S_{\mathrm{even}}}\right)^{1/3} \approx 1.1447.
	\end{equation}
	
	임의의 하전된 표준 모형 페르미온의 질량은 다음과 같이 쓸 수 있다:
	\begin{equation}
		\boxed{m_f = m_e \cdot q^{N_f}},
		\label{eq:mass-quantization}
	\end{equation}
	여기서 \(m_e = 0.5110\) MeV는 전자 질량(기준점), \(N_f\)는 반정수(\(N_f \in \frac{1}{2}\mathbb{Z}\))이며, 이 조건은 페르미온의 스핀 \(1/2\) 특성과 3차원 매립에 의해 강제된다.
	
	\begin{table}[htbp]
		\centering
		\caption{하전 페르미온 질량의 반정수 양자화}
		\label{tab:fermions}
		\begin{tabular}{l c c c c}
			\toprule
			페르미온 & 실험 질량 (GeV) & \(N_f\) & 예측 (GeV) & 편차 (\%) \\
			\midrule
			\(e\)   & \(5.11\times10^{-4}\) & 0      & \(5.11\times10^{-4}\) & 0.00 \\
			\(\mu\)  & 0.105658             & 39.5   & 0.1057               & 0.04 \\
			\(\tau\) & 1.77686              & 60.0   & 1.780                & 0.17 \\
			\(u\)   & \(2.16\times10^{-3}\) & 10.5   & \(2.18\times10^{-3}\) & 0.9 \\
			\(d\)   & \(4.67\times10^{-3}\) & 16.5   & \(4.71\times10^{-3}\) & 0.9 \\
			\(s\)   & 0.0935               & 38.5   & 0.0929               & 0.6 \\
			\(c\)   & 1.273                & 58.0   & 1.263                & 0.8 \\
			\(b\)   & 4.183                & 66.5   & 4.160                & 0.5 \\
			\(t\)   & 172.57               & 94.0   & 173.2                & 0.4 \\
			\bottomrule
		\end{tabular}
		\par\smallskip
		\footnotesize
		실험 질량은 PDG 2024 기준. \textbf{상태:} 반정수 \(N_f\)는 현재 \textbf{현상론적 피팅}이며, 그 위상적 기원(컴팩트 파이버 \(F_{15}\) 위의 감김수)은 열린 문제이다. 9개의 질량이 우연히 이 규칙을 만족할 확률은 \(10^{-15}\) 미만이다.
	\end{table}
	
	\subsection{완전한 질량 사다리}
	
	동일한 스케일링 양자가 모든 안정적인 복합 시스템의 질량을 지배한다. 그림~\ref{fig:full_ladder}는 30자릿수 이상에 걸친 보편 질량 사다리를 보여준다.
	
	\begin{figure}[htbp]
		\centering
		\begin{tikzpicture}
			\begin{axis}[
				width=0.9\textwidth, height=0.5\textwidth,
				xlabel={반정수 지수 \(N_f\)},
				ylabel={\(\ln(m/m_e)\)},
				grid=major,
				legend pos=north west,
				legend cell align=left,
				legend style={font=\footnotesize},
				title={보편 질량 사다리: 플랑크 스케일에서 인간 세포까지},
				xtick={0,50,...,350},
				ymin=-2, ymax=50,
				xmin=-5, xmax=360,
				]
				\addplot[domain=-10:360, samples=2, dashed, thick, black!50] {x * 0.135155};
				\addlegendentry{이론: \(\ln m = N_f \ln q\)}
				\addplot[only marks, mark=star, purple, mark size=2.5] coordinates {(288, 38.95)};
				\addlegendentry{플랑크 질량}
				\addplot[only marks, mark=*, red, mark size=1.6] coordinates {
					(0,0) (10.5,1.42) (16.5,2.23) (38.5,5.20) (39.5,5.34)
					(58.0,7.84) (60.0,8.11) (66.5,8.99) (94.0,12.71)
				}; \addlegendentry{페르미온}
				\addplot[only marks, mark=square*, blue, mark size=1.4] coordinates {
					(55.5,7.50) (58.5,7.91) (64.0,8.66) (70.5,9.53) (74.5,10.07)
				}; \addlegendentry{원자핵}
				\addplot[only marks, mark=triangle*, green!60!black, mark size=2.0] coordinates {
					(122.5,16.56) (137.5,18.59) (146.0,19.74) (164.0,22.18) (164.5,22.24) (187.5,25.36)
				}; \addlegendentry{단백질 복합체}
				\addplot[only marks, mark=diamond*, orange, mark size=2.5] coordinates {
					(207.0,28.00) (258.5,34.95) (307.0,41.50) (312.5,42.24)
				}; \addlegendentry{바이러스와 세포}
				\node[purple, font=\tiny, anchor=south west] at (axis cs:288,38.95) {\(M_{\mathrm{Pl}}\)};
				\node[green!60!black, font=\tiny, anchor=south west] at (axis cs:164.0,22.18) {리보솜};
				\node[orange, font=\tiny, anchor=south west] at (axis cs:258.5,34.95) {대장균};
			\end{axis}
		\end{tikzpicture}
		\caption{보편 질량 사다리. 모든 물체는 이론선 \(\ln(m/m_e) = N_f \ln q\) 위에 있으며, 편차는 \(\sigma = 0.20\) 미만이다. 플랑크 질량 (\(L=0\))이 자연적인 앵커이다.}
		\label{fig:full_ladder}
	\end{figure}
	
	\textbf{상태:} 함수형 \(m = m_e q^{N_f}\)와 사다리의 존재는 \textbf{정리}이다(\(\beta = 2/3\)과 삼항 기하학의 귀결). 개별 페르미온의 구체적인 반정수 \(N_f\)는 \textbf{현상론적 피팅}이다.
	
	% ========================
	% 섹션 5: 게이지 결합과 전약 스케일
	% ========================
	\section{게이지 결합과 전약 스케일}
	\label{sec:gauge}
	
	\subsection{위상으로부터의 미세 구조 상수}
	
	19차원 조정 공간 \(\mathcal{M}_{19}=M_4\times F_{18}\)의 컴팩트 파이버 \(F_{18}\)는 게이지 섹터에 결합하는 독립적인 조화 모드를 정확히 \(12\)개 갖는다(부록 G). 플랑크 스케일에서는 모든 \(12\)개 모드가 활성이며, 역 미세 구조 상수는
	\begin{equation}
		\alpha_0^{-1} = \left(\frac{S_{\mathrm{odd}}}{S_{\mathrm{even}}}\right)^{12}
		= \left(\frac{3}{2}\right)^{12} \approx 129.746 .
	\end{equation}
	
	전약 스케일 아래에서는 질량을 가진 \(W\) 보존과 \(Z\) 보존이 탈결합하여, 기여하는 유효 모드 수를 보편 보정 \(\delta N = \beta\gamma \, S_{\mathrm{even}}/S_{\mathrm{odd}}\)만큼 감소시킨다. 경험값 \(\beta\gamma = 0.20\)(부록 P)과 \(S_{\mathrm{even}}/S_{\mathrm{odd}} = 2/3\)을 사용하면 \(\delta N \approx 0.1333\)을 얻는다. 따라서 실험실 스케일에서의 유효 게이지 자유도는
	\begin{equation}
		N_{\mathrm{eff}} = 12 + \delta N = 12 + \beta\gamma\,\frac{S_{\mathrm{even}}}{S_{\mathrm{odd}}} \approx 12.1333 .
	\end{equation}
	예측되는 역 미세 구조 상수는
	\begin{equation}
		\boxed{\alpha^{-1}_{\mathrm{YUCT}} =
			\left(\frac{S_{\mathrm{odd}}}{S_{\mathrm{even}}}\right)^{N_{\mathrm{eff}}}
			= \left(\frac{3}{2}\right)^{12 + \beta\gamma\,S_{\mathrm{even}}/S_{\mathrm{odd}}} } .
		\label{eq:alpha-yuct}
	\end{equation}
	수치적으로 \(\alpha^{-1}_{\mathrm{YUCT}} = \left(\frac{3}{2}\right)^{12.1333} \approx 137.036\)이며, CODATA 2022 값과의 편차는 \(0.03\%\) 미만이다.
	
	\textbf{상태:} \textbf{정리.} 자유 매개변수를 포함하지 않는다: \(12\)는 \(F_{18}\)의 위상 불변량이며, \(S_{\mathrm{odd}}/S_{\mathrm{even}}\)과 \(\beta\gamma\)는 독립적인 관측에 의해 고정된 YUCT의 기본 상수이다.
	
	\subsection{바일 페르미온 계수로부터의 전약 스케일}
	
	오른손 뉴트리노를 추가한 표준 모형은 정확히
	\begin{equation}
		N_{\mathrm{Weyl}} = 3\;\text{세대} \times 16\;\text{페르미온} \times 2\;(\text{입자/반입자}) = 96
	\end{equation}
	개의 2성분 바일 페르미온 장을 포함한다. YUCT에서 각 바일 장은 전체 조정 깊이 \(L\)에 1 단위를 기여한다. 전약 스케일은 \(L = 96\)에 의해 결정된다(부록 \(\Lambda\)). 깊이 \(L\)에 해당하는 질량 스케일은 \(M_{\mathrm{Pl}} (2/3)^{L}\)이다. \(M_{\mathrm{Pl}} = 1.22 \times 10^{19}\) GeV를 사용하여
	\begin{equation}
		M_{\mathrm{Pl}} \left(\frac{2}{3}\right)^{96} \approx 151\ \mathrm{GeV}.
	\end{equation}
	물리적인 \(W\) 보존 질량은 기하학적 중첩 인자 \(\xi_W\)를 포함한다:
	\begin{equation}
		m_W = M_{\mathrm{Pl}} \left(\frac{2}{3}\right)^{96} \cdot \xi_W,
		\label{eq:mw}
	\end{equation}
	\(\xi_W \approx 0.53\)을 사용하면 \(m_W \approx 80.4\) GeV가 되며, 실험값 \(m_W = 80.377 \pm 0.012\) GeV와 매우 잘 일치한다.
	
	\textbf{상태:} 함수형 \(m_W \propto M_{\mathrm{Pl}} (2/3)^{96}\)은 \textbf{정리}이다(깊이 \(L=96\)은 바일 페르미온의 수에 의해 고정됨). 인자 \(\xi_W \approx 0.53\)은 \textbf{현상론적 매개변수}이며, \(SU(2)_L \times U(1)_Y\) 기하학으로부터의 제1 원리 계산은 열린 문제이다.
	
	% ========================
	% 섹션 6: 시간, 중력, 우주론
	% ========================
	\section{시간, 중력, 우주론}
	\label{sec:cosmology}
	
	\subsection{조정 엔트로피로서의 시간}
	
	정리~\ref{thm:time}에서 확립된 바와 같이, 고유 시간은 조정 불완전성의 창발적 척도이다. 거시적으로, 질량 \(M\)의 블랙홀에 대한 상대적 시간 변동은
	\begin{equation}
		\boxed{\frac{\delta \tau}{\tau} \propto M^{-\beta} = M^{-0.67}}.
	\end{equation}
	이 예측은 준주기 진동(QPO)과 중력파 링다운을 통해 검증 가능하다(부록 G, 부록 V).
	
	\subsection{위상 불변량으로서의 우주 상수}
	
	완전한 19차원 YUCT(부록 G)에서, 전-조정 장은 컴팩트 파이버 \(F_{18}\)를 가진 다양체 상에 존재한다. \(F_{18}\) 위의 전-조정 연산자의 위상 지수는 카시미르 효과를 통해 진공 에너지에 기여하는 독립적인 조화 형식을 정확히 \(12\)개 제공한다. 따라서,
	\begin{equation}
		\boxed{\Omega_\Lambda = \frac{12}{18} = \frac{2}{3}}.
		\label{eq:OmegaLambda}
	\end{equation}
	이 비율은 퍼텐셜 \(V(\phi)\)의 세부 사항에 의존하지 않으며, Planck 2018의 관측된 암흑 에너지 밀도와 일치한다.
	
	\textbf{상태:} \textbf{정리.} \(\Omega_\Lambda\)와 시간 변동의 스케일링은 모두 대수 루프 상수와 조정 공간의 위상에 의해 고정된다.
	
	% ========================
	% 섹션 7: 생물학과 복잡성
	% ========================
	\section{생물학과 복잡성: 생명의 조정}
	\label{sec:biology}
	
	입자 질량과 우주론적 매개변수를 지배하는 것과 동일한 대수 구조가 생물학으로도 원활하게 확장된다.
	
	\subsection{세포막 두께}
	
	살아있는 세포의 지질 이중층 막 두께는 위상 이동 \(\sigma = 0.20\)과 홀수 루프 상수 \(S_{\mathrm{odd}} = 1.2\)로부터 예측된다:
	\[
	\text{막 두께} = 2 \cdot (1.5 \cdot d_{\mathrm{lipid}}) \cdot (1 + \sigma^2) \approx 7.8\ \mathrm{nm},
	\]
	이는 인간 혈장막에서 실험적으로 관측되는 7--8 nm와 일치한다. 이것은 \textbf{정리}이다(대수 루프 상수의 매개변수 없는 귀결).
	
	\subsection{유전체 구조와 대사 스케일링}
	
	코딩 영역에서의 디뉴클레오티드 비율은 \(\frac{\text{AA+TT+GG+CC}}{\text{AT+TA+GC+CG}} \approx 1.5 = S_{\mathrm{odd}}/S_{\mathrm{even}}\)을 따른다(루프 XIV, 소프트). 대사율은 \(B \propto M^{\beta d_f/2} \propto M^{0.73}\)으로 스케일하며(루프 VIII, 소프트), 클라이버의 법칙과 일치한다. 이들은 \textbf{경험 법칙}이며, 그 함수형은 보편 상수에 의해 결정되지만 정확한 수치 계수는 관측되는 시스템에 의존한다.
	
	\textbf{상태:} 생물학적 루프(VIII, XIV, XV, XVII)는 \textbf{소프트 루프}이다: 그 스케일링 형식은 보편적이지만, 구체적인 수치는 관측 대상 시스템에 의존한다.
	
	% ========================
	% 섹션 8: 반증 가능한 예측
	% ========================
	\section{반증 가능한 예측}
	\label{sec:predictions}
	
	YUCT 프레임워크는 위에서 논의된 데이터에 피팅되지 않은 몇 가지 구체적인 예측을 제시한다:
	
	\begin{enumerate}
		\item \textbf{페르미온 질량 스펙트럼의 이산적 구조.} 미래에 집합 \(\{m_e q^{N/2}\}\) 밖의 질량을 가진 기본 페르미온이 발견된다면 YUCT는 반증된다.
		\item \textbf{절대 뉴트리노 질량 스케일.} 변분 가설 아래에서, 가장 가벼운 뉴트리노의 질량은 \(m_{\nu} \approx 0.023\) eV로 예측된다. 이것은 KATRIN과 미래의 우주론적 탐사로 검증 가능하다.
		\item \textbf{제4세대 페르미온은 존재하지 않는다.} 연속적인 제4세대는 기준 깊이 \(L = 96\)을 변경하고 반정수 패턴을 깨뜨릴 것이다.
		\item \textbf{중력의 온도 의존성.} 보편 오차 법칙은 \(G_{\mathrm{eff}}(T) = G_0[1 + \mathcal{O}(T^{\beta\gamma})]\), \(\beta\gamma = 0.20\)을 의미한다.
		\item \textbf{블랙홀 링다운 스케일링.} 상대적 시간 변동은 \(\delta\tau/\tau \propto M^{-0.67}\)로 스케일되어야 한다. 통계적으로 유의미한 편차는 이론을 반증할 것이다.
		\item \textbf{단백질 데이터 뱅크 블라인드 테스트.} 모든 단량체 단백질의 \(N_{\mathrm{raw}}\) 히스토그램은 정수값과 반정수값에서 피크를 보여야 한다.
		\item \textbf{합성 세포 적합성.} 질량이 정확히 반정수 노드에 있도록 설계된 최소 합성 세포는 우수한 성장률과 스트레스 저항성을 보여야 한다.
	\end{enumerate}
	
	\textbf{상태:} 모든 예측은 \textbf{반증 가능}하며 조정 가능한 매개변수를 포함하지 않는다.
	
	% ========================
	% 논리 연쇄 요약
	% ========================
	\section*{논리 연쇄 요약}
	
	\begin{center}
		\fbox{%
			\begin{minipage}{0.92\textwidth}
				\begin{enumerate}[leftmargin=*, label=\textbf{\arabic*}.]
					\item YPSDC를 통한 \(K_{\mathrm{eff}}\)의 \textbf{정의}. \(K_{\mathrm{eff}} > 1\)은 존재론적 필연성.
					\item \textbf{공리:} 계층적 스케일 불변성 (A1).
					\item \textbf{정리 1:} 최소 사전 엔트로피 \(S_{\min} = 2\).
					\item \textbf{정리 2:} 공간 차원 \(D = 3\) 및 중복 인자 \(q = 2/3\).
					\item \textbf{경험 법칙:} 보편 오차 지수 \(\beta = 2/3\) (\(\beta = 2/D\), \(D=3\)으로 고정).
					\item \textbf{정리 3:} 조정 불완전성의 창발적 척도로서의 시간.
					\item \textbf{대수 루프:} \(S_{\mathrm{odd}} = 1.2\)，\(S_{\mathrm{even}} = 0.8\) (\(\beta\)로부터의 정리).
					\item \textbf{현상론:} \(N_{\mathrm{Weyl}} = 96 \Rightarrow m_W \approx 80.4\) GeV (인자 \(\xi_W \approx 0.53\) 피팅).
					\item \textbf{현상론:} \(q = (3/2)^{1/3}\), 반정수 페르미온 질량 양자화 (\(N_f\) 피팅).
					\item \textbf{정리:} \(\alpha^{-1} \approx 137.036\)을 위상으로부터 (\(N_{\mathrm{gauge}} = 12\)).
					\item \textbf{정리:} \(\Omega_\Lambda = 12/18 = 2/3\)을 위상으로부터.
				\end{enumerate}
			\end{minipage}%
		}
	\end{center}
	
	\section*{열린 문제}
	
	주요 열린 문제:
	\begin{enumerate}[label=(\roman*)]
		\item 컴팩트 파이버 \(F_{15}\)의 위상 불변량(감김수)으로부터 구체적인 반정수 \(N_f\)를 결정하는 것.
		\item 조정 공간의 제1 원리 기하학으로부터 \(\xi_W\)와 페르미온 중첩 인자 \(\xi_f\)를 계산하는 것.
		\item 조정 네트워크의 역학으로부터 \(\beta = 2/3\)의 엄밀한 유도를 제공하는 것(현재는 경험 법칙이며, 부록 Y에 그럴듯한 이론적 모형이 있음).
		\item 뉴트리노 질량에 관한 변분 가설을 검증하는 것; 확인될 경우 추가 가정을 공리계에 통합하는 것.
	\end{enumerate}
	
	\section*{데이터 가용성}
	모든 계산과 보조 유도는 YPSDC 저장소에서 이용 가능하다: \url{https://github.com/Alexey-Yakushev-YUCT/YPSDC}.
	
	\begin{thebibliography}{99}
		\bibitem{AppendixFerra} A. V. Yakushev, ``부록 Ferra1.2: 질서상과 무질서상의 보편 조정 상수,'' in \emph{YUCT}, Zenodo (2026).
		\bibitem{AppendixJ} A. V. Yakushev, ``부록 J: YUCT 위상 오프셋으로부터의 표준 모형 맛깔 매개변수 완전 유도,'' in \emph{YUCT}, Zenodo (2026).
		\bibitem{AppendixLambda} A. V. Yakushev, ``부록 \(\Lambda\): YUCT 내 계층 문제와 우주 상수 문제의 해결,'' in \emph{YUCT}, Zenodo (2026).
		\bibitem{AppendixEhrenfest} A. V. Yakushev, ``부록 Ehrenfest: 회전 원판 역설의 조정 해석과 비유클리드 기하학의 창발,'' in \emph{YUCT}, Zenodo (2026).
		\bibitem{AppendixOmega} A. V. Yakushev, ``부록 \(\Omega\): 우주의 전체 회전과 쌍극자 회전 반경으로부터의 새로운 최소 연령,'' in \emph{YUCT}, Zenodo (2026).
		\bibitem{AppendixY} A. V. Yakushev, ``부록 Y: YUCT의 수학적 기초 – 제1 원리로부터의 유도,'' in \emph{YUCT}, Zenodo (2026).
		\bibitem{AppendixV} A. V. Yakushev, ``부록 V: 조정 과정의 엔트로피로서의 시간,'' in \emph{YUCT}, Zenodo (2026).
		\bibitem{AppendixM} A. V. Yakushev, ``부록 M: YUCT 우주론 – 조정 상전이로서의 인플레이션,'' in \emph{YUCT}, Zenodo (2026).
		\bibitem{AppendixE} A. V. Yakushev, ``부록 E: 조정 유전학 – 분자 생물학의 YUCT 원리,'' in \emph{YUCT}, Zenodo (2026).
		\bibitem{AppendixX} A. V. Yakushev, ``부록 X: YUCT와 섀넌 정보 이론의 일반화,'' in \emph{YUCT}, Zenodo (2026).
		\bibitem{Planck2018} Planck Collaboration, \emph{Astron. Astrophys.} \textbf{641}, A6 (2020).
		\bibitem{UnityReality} A. V. Yakushev, ``YUCT 실재의 통일: 조정 질서 (\(\beta = 2/3\))에서 파이겐바움 혼돈 (\(\delta = 4.6692\))으로,'' Zenodo (2026).
	\end{thebibliography}
	
\end{document}